\section{The structural limit passages}\label{sec:stability}

\subsection{Linear identities and approximate symmetries}
Suppose \(P_j\Rightarrow P\) on \(H_{-t}\) and
\begin{equation}\label{eq:quadratic}
 \E_{P_j}|F(f)|^2\le C\|f\|_2^2
 \qquad(f\in\mathcal S\text{ real})
\end{equation}
with one finite \(C\).
Under \cref{hyp:identity}, exact annihilation of \(v_j(\xi)\) gives
\begin{equation}\label{eq:identity-residual}
 \E_{P_j}|F(D^*\xi)|^2
 \le C\|D^*\xi-v_j(\xi)\|_2^2\longrightarrow0.
\end{equation}
For \(M>0\), the function
\(\min\{|F(D^*\xi)|^2,M\}\) is bounded and continuous on \(H_{-t}\).
Its limiting expectation is zero, so \(F(D^*\xi)=0\),
\(P\)-almost surely. Choose a countable dense family of smooth tests
on each member of a compact exhaustion of \(\R^d\), and intersect
the corresponding full-measure events. Continuity of \(D^*\) and
of each sample distribution extends the identity to every compact
smooth test.

For symmetries, put \(\chi_j(f)=\E_{P_j}e^{iF(f)}\).
The inequality \(|e^{iu}-e^{iv}|\le|u-v|\) and
\eqref{eq:quadratic} imply
\begin{equation}\label{eq:character-continuity}
 |\chi_j(f)-\chi_j(h)|\le\sqrt C\,\|f-h\|_2
 \qquad(f,h\in\mathcal S\text{ real}).
\end{equation}
For a fixed \(g\) in \cref{hyp:symmetry}, this gives
\[
 |\chi_j(g\cdot f)-\chi_j(f)|
 \le\sqrt C\,\|g\cdot f-T_{j,g}f\|_2\longrightarrow0.
\]
Weak convergence passes both fixed-test characteristic functions
to the limit. Applying the equality to all real linear combinations
of any finite family proves equality of its joint laws.
These evaluations generate the sigma-algebra, so \(P\) is
\(g\)-invariant.

In particular, suppose \(P_j\) is invariant under translations in
\(a_j\Z^d\), \(a_j\to0\). Given \(b\), choose \(b_j\in a_j\Z^d\)
with \(b_j\to b\). The maps
\(T_{j,\tau_b}f=\tau_{b_j}\cdot f\) meet the approximation
hypothesis by continuity of translations on \(L^2\).
Exact fixed symmetries use \(T_{j,g}f=g\cdot f\).

\subsection{Reflection forms}
For any law of tempered distributions, the linear span of
\(e^{iF(f)}\), \(f\in C_c^\infty(O;V)\) real, is dense in
\(L^2(\Sigma_O)\). To see this, let \(B\) be orthogonal to these
characters and condition it on finitely many evaluations.
Orthogonality to every real linear combination of their tests
says that the Fourier transform of the resulting finite complex
measure is zero. Fourier uniqueness makes that conditional
expectation zero. A countable dense family of compact tests
generates \(\Sigma_O\), since every sample is a continuous
distribution. Conditional expectations on its increasing finite
evaluation sigma-algebras converge in \(L^2\), proving the claim.

Fix \(r\in\mathcal R\). For
\(B(F)=\sum_bz_be^{iF(f_b)}\), the reflected form equals
the left side of \eqref{eq:rp-approx}. Each summand is bounded
and continuous on \(H_{-t}\), so weak convergence gives
\[
 \E_P[\overline{B(r\cdot F)}B(F)]\ge0.
\]
The preceding symmetry argument gives \(r\)-invariance of \(P\).
Consequently
\begin{equation}\label{eq:reflection-form-bound}
 \big|\E_P[\overline{B(r\cdot F)}C(F)]\big|
 \le\|B\|_{L^2(P)}\|C\|_{L^2(P)}.
\end{equation}
Density extends positivity to all
\(B\in L^2(\Sigma_{H_r^+},P)\).
If the plane of \(r\) is \(n\cdot x=0\), its translate at height
\(c\) has reflection
\[
 r_c=\tau_vr\tau_{-v},\qquad v=c|n|^{-2}n.
\]
Translation invariance identifies the two forms and carries their
positive-side sigma-algebras to one another. It also gives
\(r_c\)-invariance.

\begin{proof}[Proof of \cref{thm:structural-limit}]
The exponential bound implies
\(\E_{P_j}F(f)=0\) and
\(\E_{P_j}|F(f)|^2\le\|f\|_2^2\);
see \cref{prop:subgaussian}.
The preceding arguments with \(C=1\) prove the identities,
invariance and reflection positivity.
The Gaussian bound for \(P\) and convergence of all finite-family
moments and cumulants follow from
\cref{thm:gaussian-convergence}, whose hypotheses are
the weak convergence and \eqref{eq:common-gaussian}.
\end{proof}

